% =============================================================================
% Ethics Statement and Broader Impact
% Self-contained file: % =============================================================================
% Ethics Statement and Broader Impact
% Self-contained file: % =============================================================================
% Ethics Statement and Broader Impact
% Self-contained file: % =============================================================================
% Ethics Statement and Broader Impact
% Self-contained file: \input{ethics_statement} into main.tex
% Intended placement: after the Conclusion section, before References.
% Label: \label{sec:impact_standalone}
% Target length: >= 2 pages (NeurIPS Broader Impact guidance).
% =============================================================================

\section{Ethics, Limitations, and Broader Impact}
\label{sec:impact_standalone}

This statement is written to satisfy the NeurIPS Code of Ethics and Broader
Impact requirements. It complements the shorter Limitations and Broader
Impact subsections embedded in Section~\ref{sec:limitations} by consolidating
in one place a dual-use analysis, unresolved compute-accounting discrepancies, carbon
scenarios, data provenance disclosures, a candid acknowledgment
of the closed-source training backend on which a fraction of our headline
numbers depends, and a list of methodological limits we are aware of but
did not have resources to close within the submission window.

% ---------------------------------------------------------------------------
\subsection{Dual-Use Analysis: Misuse Risks of Reasoning RL}
\label{sec:dualuse}
% ---------------------------------------------------------------------------

\paragraph{Threat model.}
\ours{} releases (i) training scripts that apply GRPO to the GSM8K
\cite{cobbe2021training} mathematical reasoning benchmark and to the
Salesforce xLAM-Function-Calling-60k \cite{liu2024xlam} tool-use corpus,
(ii) a set of LoRA adapters published on the HuggingFace Hub
(\texttt{arvindcr4/tinker-rl-bench-*}), and (iii) diagnostic code for
Zero-Variance Fraction, reward-stability, and length-bias analysis. We
analyse four concrete misuse pathways these artefacts could, in principle,
enable, and we describe the existing guardrails and the residual risk we
judge acceptable for publication.

\paragraph{M1. Weaponisation of mathematical reasoning.}
The most capability-relevant artefact we release is GRPO training code evaluated
on grade-school arithmetic and multi-step reasoning; the reviewed evidence is
inconclusive about a capability improvement. A malicious
operator could attempt to extend this pipeline to tasks of greater
dual-use concern, e.g., chemistry olympiad problems, cryptographic key
recovery, or financial fraud. We judge the \emph{marginal} uplift provided
by our release to be small for three reasons. First, GRPO at our training
scale does not create new capabilities; the pre-RL held-out accuracy of
Qwen3-8B-Instruct for GSM8K (Section~\ref{sec:results}, 82.0\%) shows that the
bulk of held-out accuracy is attributable to the checkpoint's pretraining and
instruction tuning. Our 5-seed one-sample seed-level comparison against a fixed
baseline estimates \(+1.3\) percentage points ($p{=}0.26$); it is neither
paired-seed nor reviewed-record evidence of improvement, and should not be read
as a bound on capability change: the capability difference is unmeasured within
wide uncertainty (single base pass, no baseline variance), so ``inconclusive''
does not establish ``small uplift.''
Second, the public literature already contains GRPO implementations
\cite{shao2024deepseekmath, openrlhf2024, verl2024, vonwerra2022trl} that are at
least as capable. Third, any reasoning uplift is bounded by the rewarder:
GSM8K rewards use exact-match against a human-written gold answer, which
does not generalise to the domains listed above without a new reward
signal, which itself requires expertise and labelled data.

\paragraph{M2. Reward hacking and the long-term alignment risk.}
GRPO, like all policy-gradient RL from a proxy reward, is vulnerable to
reward hacking: the model learns to exploit idiosyncrasies of the reward
function rather than solve the intended task \cite{skalse2022specificationgaming,
krakovna2020specification}. Our length-bias discussion
(Section~\ref{sec:length_bias}) records suspicious reward trajectories, not
a verified causal demonstration of reward hacking.
Users who apply our scripts to under-specified reward functions in
production settings (e.g., a custom ``helpfulness'' classifier) may observe
models that appear to improve on held-out metrics while silently
optimising against unintended proxies. We include the Zero-Variance
Fraction and reward-stability diagnostics precisely as retrospective summaries; they do not identify policy drift or reward hacking.

\paragraph{M3. Circumventing safety training via distillation or tool-use.}
The tool-use track trains LoRA adapters that teach a base model to emit
schema-valid JSON function calls \cite{liu2024xlam}. A motivated adversary
could in principle use the same pipeline to teach a base model to emit
jailbreak prompts as ``tool calls'' or to invoke an adversarial external
tool. We mitigate this by (a) releasing only adapters, not merged weights,
while noting that an adapter can still alter safety behavior and does not
guarantee that instruct-model safeguards remain effective; (b) documenting in model
cards that the tool-use adapters are trained on a narrow schema (five
synthetic tools) and have not been safety-evaluated for open-ended function
calling; and (c) not releasing any function-calling harness that would
actually execute emitted calls.

\paragraph{M4. Compute-efficiency externalities.}
This audit makes no demonstrated efficiency claim or group-size recommendation. Improved training efficiency is
itself dual-use: it lowers the cost of running adversarial fine-tuning at
scale. We judge this externality to be modest relative to the already-low
cost of fine-tuning a 0.6B--8B model with commodity LoRA recipes
\cite{hu2022lora}, but we flag it explicitly.

\paragraph{Residual risk.}
After weighing M1--M4 we judge that the reproducibility, calibration, and
pedagogical benefits of releasing \ours{} outweigh the residual misuse
risk. We adopt the standard mitigation of releasing adapters (not merged
weights), documenting intended use and out-of-scope use in model cards,
and publishing alongside the diagnostics a researcher would need to detect
reward hacking.

% ---------------------------------------------------------------------------
\subsection{Compute Cost Accounting}
\label{sec:compute_accounting}
% ---------------------------------------------------------------------------

Table~\ref{tab:compute_spend} retains historical platform entries for audit,
not verified project accounting. Excluding Colab, their numeric ranges sum to
USD 147--177, inconsistent with the earlier USD 130--140 total. Colab is an
additional per-person charge with an unspecified count. The main-text subset
also previously differed. Cash versus credits, dates, cohort overlap, and
failed-run inclusion have not been reconciled; no project spend total is certified.

\begin{table}[ht]
\centering
\small
\caption{Historical compute-spend entries; unreconciled cohorts and cash/credit status.
The sum is an arithmetic check, not verified expenditure \cite{henderson2018deep}.}
\label{tab:compute_spend}
\begin{tabular}{llrr}
\toprule
\textbf{Platform} & \textbf{Use} & \textbf{Runs} & \textbf{Cost (USD)} \\
\midrule
Tinker SDK v0.16.1 & GRPO on 4B/8B (original suite)        & 25 & \$40--55 \\
Tinker SDK v0.16.1 & 10x Structural Ceiling ablation       & 32 & \$65 \\
Tinker SDK v0.16.1 & World-Class Suite (frontier/MoE)      & 13 & \$25--30 \\
Modal H100         & PPO baselines (Qwen3-8B, Llama-3.1-8B)& 2  & \$12--18 \\
Modal H100         & KL-tracking run (failed)              & 1  & \$2--4   \\
Google Colab Pro (T4) & 0.5B--3B QLoRA SFT+GRPO            & --- & \$10/person \\
NVIDIA L4 (TRL baseline)& Qwen2.5-0.5B GSM8K, 5 seeds      & 5  & \$3--5   \\
HuggingFace Hub    & Model + adapter hosting               & --- & \$0 \\
Weights \& Biases  & Experiment tracking (academic tier)   & --- & \$0 \\
\midrule
\multicolumn{3}{l}{\textbf{Project total (unreconciled)}} & \textbf{not available} \\
\bottomrule
\end{tabular}
\end{table}

\paragraph{Hidden costs.} Beyond direct platform spend, the project
consumed human labour (approximately six student-FTE-months, unpaid) and
infrastructure generously subsidised by PES University's LLM Research
Group (shared access to one A100 node used for TRL baselines and
pre-submission sanity checks). No author received commercial compensation
from Thinking Machines, Modal, or any other vendor mentioned.

\paragraph{Failed-run accounting.} Of the Tinker spend, we estimate
\textasciitilde\$30--35 was consumed by runs that ultimately failed
(JWT expiry, API stalls for GPT-OSS-20b and Kimi-K2 before re-run,
KL-tracking gradient bug). We disclose this because reviewers often ask
whether reported total spend reflects only clean, successful runs; it does
not. The failed-run subset is not sufficiently reconciled to estimate a verified
percentage of total spend.

% ---------------------------------------------------------------------------
\subsection{Carbon Footprint Estimation}
\label{sec:carbon}
% ---------------------------------------------------------------------------

Table~\ref{tab:carbon2} preserves conditional energy/carbon scenarios
from historical assumptions, not measured emissions or a project total. We follow the methodology of
\citet{patterson2024empirical} and \citet{strubell2019energy}: for each
GPU-hour we multiply device TDP by wall-clock GPU-hours and the data-centre
Power Usage Effectiveness (PUE), then multiply by the grid carbon intensity
in the region where the hardware is believed to be located.

\paragraph{Assumed constants.}
\begin{itemize}
  \item NVIDIA H100 SXM5 TDP: 700\,W.
  \item NVIDIA A100 80GB TDP: 400\,W.
  \item NVIDIA L4 TDP: 72\,W.
  \item NVIDIA T4 TDP: 70\,W.
  \item Data-centre PUE: 1.1 (conservative; typical hyperscale PUE is
        1.10--1.15 \cite{google2023pue, aws2023pue}).
  \item US average grid intensity (EPA eGRID 2023): 0.367\,kg\,CO$_2$-eq / kWh
        \cite{epa2023egrid}.
  \item Indian average grid intensity (CEA 2023): 0.716\,kg\,CO$_2$-eq / kWh
        \cite{cea2023co2}.
  \item We use the US average for Modal H100 (US-East region as configured)
        and Tinker (location undisclosed; assumed US), and the Indian
        average for Colab Pro T4 used by authors based in Bengaluru
        (Google Cloud asia-south1).
\end{itemize}

\paragraph{GPU-hour estimates.}
Tinker GPU-hour counts are not disclosed by the platform. We estimate them
from wall-clock run duration and assume a single H100-class accelerator per
job, which is an unverified hardware-allocation assumption. The 18 Modal GPU-hours
in this scenario conflict with the 250 GPU-hours in Appendix~\ref{app:compute};
neither is selected as correct without a cohort-specific billing ledger.

\begin{table}[ht]
\centering
\small
\caption{Historical conditional carbon scenarios, not measured emissions. Tinker hardware
is undisclosed; 1$\times$ H100 per run is an assumption. Cohort boundaries and
failed-run coverage are unresolved; the scenario total is not a project estimate.}
\label{tab:carbon2}
\begin{tabular}{lrrrrr}
\toprule
\textbf{Platform} & \textbf{GPU-h} & \textbf{TDP\,(W)} & \textbf{PUE} &
\textbf{Energy\,(kWh)} & \textbf{CO$_2$\,(kg)} \\
\midrule
Tinker (assumed H100)   & \textasciitilde 950 & 700 & 1.1 & \textasciitilde 732 & \textasciitilde 269 \\
Modal H100 (US-East)    & \textasciitilde 18  & 700 & 1.1 & \textasciitilde 14  & \textasciitilde 5.1 \\
PES A100 (in-kind)      & \textasciitilde 60  & 400 & 1.1 & \textasciitilde 26  & \textasciitilde 19   \\
Colab Pro T4 (asia-south1) & \textasciitilde 40 & 70 & 1.2 & \textasciitilde 3.4 & \textasciitilde 2.4  \\
NVIDIA L4 (TRL baseline)   & \textasciitilde 10 & 72 & 1.1 & \textasciitilde 0.8 & \textasciitilde 0.3  \\
\midrule
\textbf{Scenario sum (not project total)} & & & & \textasciitilde 776 & \textasciitilde 296 \\
\bottomrule
\end{tabular}
\end{table}

\paragraph{Interpretation and uncertainty.}
The tabulated sum is conditional on unreconciled hours, assumed accelerator
counts, TDP, PUE, and locations. It is not a verified footprint or an upper bound:
unknown allocation can raise or lower actual consumption. Omitting failed jobs
in a future replication does not establish a bound on this project. A release
needs a single ledger with dates, cohort IDs, GPU counts, measured versus assumed
hours, credits versus cash, and failure status before reporting totals.

\paragraph{Offset and mitigation.}
We did not purchase carbon offsets. Instead, we have (a) made available selected
checkpoint artifacts so that downstream users do not
need to re-run our experiments; (b) retained local step-level CSV logs so that
learning curves can be inspected without re-training; and (c) documented
in Section~\ref{sec:compute_accounting} which runs were wasteful, so that
replicators can skip them. We regard reproducibility-by-artefact as a
more durable mitigation than offsets.

% ---------------------------------------------------------------------------
\subsection{Data Provenance}
\label{sec:provenance}
% ---------------------------------------------------------------------------

\ours{} uses only publicly released research datasets. No private data,
personally identifiable information, or licensed proprietary content is
used. No human annotators were employed during this work. We document
each dataset below.

\paragraph{GSM8K \cite{cobbe2021training}.}
Grade-school math word problems (7{,}473 train / 1{,}319 test)
authored by human writers contracted by OpenAI. Released under the
\textbf{MIT License} via \url{https://github.com/openai/grade-school-math}.
We use the standard train/test splits without modification. The canonical
citation is \citet{cobbe2021training}. Known limitations of GSM8K include
gender-neutral but culturally US-centric word problems (names, currencies,
sports) and a moderate rate of human-labelling errors estimated at
\textasciitilde 2\% by \cite{zhang2024gsm1k}. Our held-out evaluation
respects the test/train split.

\paragraph{Salesforce xLAM-Function-Calling-60k \cite{liu2024xlam}.}
60{,}000 function-calling examples generated by Salesforce Research as
training data for the xLAM agentic model family. Released under the
\textbf{CC BY 4.0} license via \url{https://huggingface.co/datasets/Salesforce/xlam-function-calling-60k}.
We use the public release as-is; specifically, we use the first 35 prompts
$\times$ 10 rollouts in the 10x Structural Ceiling tool-use track and the
full 60k split for the xlam-60k real-data run in Section~\ref{sec:results}.
Known limitations: xLAM schemas are synthetic and skew toward cleanly-typed
arguments; the dataset under-represents ambiguous, error-handling, and
multi-turn tool-calling. Attribution to Salesforce is required by the
license and is provided in Section~\ref{sec:acknowledgments} and in the
xLAM-derived model cards.

\paragraph{HumanEval \cite{chen2021codex}.}
164 Python programming problems with unit tests, released by OpenAI under
the \textbf{MIT License}. The local restricted sandbox does not implement an unmodified standard
pass@$k$ evaluation; code results are harness-bound reward probes. Known
limitations: small scale, English-language docstrings, single-language
coverage; documented contamination concerns against frontier pretraining
corpora \cite{riddell2024contamination}.

\paragraph{NuminaMath \cite{li2024numinamath}.}
Used by collaborator Mohammad Rafi for the multi-stage
GSM8K+NuminaMath pipeline. Released under \textbf{Apache 2.0}. We use a
downstream checkpoint reported by Rafi; our repository contains only his
scripts and the Apache-licensed derivative weights.

\paragraph{Open-Platypus \cite{lee2023platypus}.}
3{,}000 SFT examples used by collaborator Madhu for Qwen3-8B code
generation warm-up. Open-Platypus is a curated subset released under
\textbf{CC BY-NC 4.0}; the non-commercial restriction is respected in our
release, which is academic-use only.

\paragraph{Synthetic tool-use corpus (authored).}
We additionally generated a small five-tool synthetic corpus (calculator,
web-search stub, calendar, file-read, email-send) used for the reviewed descriptive
all-wrong observations. The corpus is authored by the paper's authors from
publicly documented APIs, contains no third-party content, and is
released under the \textbf{MIT License} in our repository under
\texttt{data/synthetic\_tools/}. Generation prompts are included for
auditability. No user data, scraped web content, or proprietary API
traces are present.

\paragraph{No web scraping, no human subjects.}
We did not scrape any website. We did not run any study with human
participants. No IRB review was therefore required. No data that could
reasonably be considered to contain PII, copyrighted content, or sensitive
personal attributes was used at any stage of training, evaluation, or
reward modelling.

% ---------------------------------------------------------------------------
\subsection{Closed-Source Tinker Acknowledgment}
\label{sec:tinker_closed}
% ---------------------------------------------------------------------------

A central tension in \ours{} is that a substantial fraction of our
headline numbers come from a closed-source commercial platform while our
paper simultaneously advocates for reproducibility and platform
independence. We do not resolve this tension by hiding it; we describe
it precisely.

\paragraph{What Tinker is and is not.}
Tinker \cite{thinkingmachines2024tinker} is a managed LLM fine-tuning and
inference service provided by Thinking Machines, Inc. It exposes a Python
SDK (\texttt{tinker==0.16.1}) that accepts custom loss functions
(\texttt{forward\_backward\_custom}), a limited set of optimisers, and
standard LoRA hyperparameters. Local scripts specify reward normalization, baseline subtraction, custom loss,
and requested LoRA/optimizer configuration. Undisclosed server-side details
include hardware placement, scheduling, execution numerics, and some batching
and sampler behavior. These are separate from the inspectable local objective.

\paragraph{What this means for our claims.}
We have not ruled out seed variance, checkpoint/model-size confounding, evaluator
differences, or backend effects. The arithmetic/GSM8K cross-library identity is
unresolved. Tinker results do not establish an upper bound or algorithmic
superiority. verl/OpenRLHF Task~4 dry runs are engineering artifacts, not
measured alternatives. The contribution remains the centered-reward identity,
descriptive audited cells, and reporting contract.

\paragraph{Reproducibility commitments we can make.}
Section~\ref{sec:reproducibility} distinguishes local source, observed
configuration, undisclosed internals, and evidence access. A dagger denotes only
interrupted execution; backend provenance is stated explicitly. Exact checkpoint
revisions and public access must be verified per artifact; no blanket release
or future rerun commitment is made here.

\paragraph{Key rotation and credential hygiene.}
Our Tinker API key (prefix \texttt{tml-...}) is held in a repository
\texttt{.env.example} placeholder only. The real key is stored in the
authors' password manager and rotated whenever a failure mode suggests
possible exfiltration. No Tinker key is committed to git history in
either the primary repository (\texttt{pes-llm-research/tinker-rl-lab}) or
the mirror (\texttt{arvindcr4/tinker-rl-lab}).

% ---------------------------------------------------------------------------
\subsection{Known Methodological Limits}
\label{sec:methodlimits}
% ---------------------------------------------------------------------------

In addition to the infrastructure failures and platform-specific caveats
enumerated in Section~\ref{sec:limitations}, we flag the following
methodological limits.

\paragraph{Short training horizons.}
Most reviewed Tinker runs have short, often 30-step horizons; this is not
a universal limit for all later scripts or campaigns. This is a
deliberate cost-control choice and means our Tinker numbers are
\emph{early-training snapshots}. We cannot rule out that longer training
would change the qualitative picture. The later 82\%$\to$83.3\% comparison
is one-sample fixed-baseline evidence, not paired-seed or reviewed-record
evidence; it is not a held-out gain whose size can be extrapolated.

\paragraph{Single-seed Tinker experiments.}
Each Tinker configuration was run once, not 3--5 times as best practice
prescribes \cite{henderson2018deep}. Tinker results therefore lack
variance estimates and do not support significance testing; we report
them descriptively. Seed-level inference requires the actual independent seed observations and
an identified estimand; reconstructed seed vectors are not eligible.

\paragraph{Train-set reward as primary Tinker metric.}
Tinker runs primarily report reward on training prompts. Only the GSM8K
track was followed up with a separate 200-example held-out evaluation
per seed. Other Tinker tracks (tool-use, xLAM) remain train-set-only and
we cannot distinguish memorisation from generalisation for those
settings.

\paragraph{LoRA only; no full fine-tuning.}
The reviewed LLM fine-tuning runs use LoRA \cite{hu2022lora}, with settings
varying by cohort; the Wave~6 rank sweep also includes rank 4. We have
\emph{not} tested whether full fine-tuning would re-order the
library/algorithm comparisons. The LoRA constraint is practical (cost) but
it means our claims about PPO vs.\ GRPO and TRL vs.\ Tinker are
LoRA-conditional.

\paragraph{Benchmark coverage is narrow.}
We evaluate four domains: GSM8K, MATH-500 (exploratory), HumanEval
(subset), and synthetic/xLAM tool-use. We do not evaluate on MT-Bench,
ArenaHard, safety benchmarks (HarmBench, ToxicChat), or truthfulness
benchmarks (TruthfulQA). Any claim about ``reasoning'' in the paper is
strictly a claim about the covered benchmarks.

\paragraph{No human preference data.}
We use verifiable rewards only (exact-match for math, unit-test for code,
schema-validity for tool-use). We do not study RLHF with human preference
data; findings should not be extrapolated to reward-model--based alignment
without further work.

\paragraph{Closed-source implementation opacity (reiterated).}
Comparisons between Tinker GRPO and TRL GRPO are confounded by the
closed-source nature of Tinker's implementation. See
Section~\ref{sec:tinker_closed} for detail.

\paragraph{Cross-platform hardware confounding.}
Tinker, Modal, Colab, and the PES A100 node use different accelerators,
different memory hierarchies, and different inference/training stacks
(Tinker proprietary, Modal vLLM, Colab PyTorch-native, TRL+Accelerate).
Observed differences in reward dynamics are not purely algorithmic.

\paragraph{Carbon estimates are approximate.}
As discussed in Section~\ref{sec:carbon}, Tinker GPU-hour and TDP are
inferred, not measured. Grid intensity is region-average, not marginal.
Numbers in Table~\ref{tab:carbon2} are conditional scenarios with
unreconciled cohorts, not measured emissions or upper bounds.

\paragraph{Exploratory, not confirmatory.}
This is an exploratory case study, not a pre-registered confirmatory
experiment. We did not pre-register primary hypotheses. The conflicting historical correction families and step-based significance
decisions are withdrawn (Section~\ref{sec:statistical-policy}); exploratory
coefficients do not establish calibrated inference.

\paragraph{Geographic and demographic limits.}
All authors are based at a single institution in Bengaluru, India. Our
choice of benchmarks, prompt phrasings, and evaluation design reflects
that context. Independent replication by groups in other regions, on
non-English tasks, and with non-Qwen / non-Llama families is an open
question.

% =============================================================================
% End of ethics_statement.tex
 into main.tex
% Intended placement: after the Conclusion section, before References.
% Label: \label{sec:impact_standalone}
% Target length: >= 2 pages (NeurIPS Broader Impact guidance).
% =============================================================================

\section{Ethics, Limitations, and Broader Impact}
\label{sec:impact_standalone}

This statement is written to satisfy the NeurIPS Code of Ethics and Broader
Impact requirements. It complements the shorter Limitations and Broader
Impact subsections embedded in Section~\ref{sec:limitations} by consolidating
in one place a dual-use analysis, unresolved compute-accounting discrepancies, carbon
scenarios, data provenance disclosures, a candid acknowledgment
of the closed-source training backend on which a fraction of our headline
numbers depends, and a list of methodological limits we are aware of but
did not have resources to close within the submission window.

% ---------------------------------------------------------------------------
\subsection{Dual-Use Analysis: Misuse Risks of Reasoning RL}
\label{sec:dualuse}
% ---------------------------------------------------------------------------

\paragraph{Threat model.}
\ours{} releases (i) training scripts that apply GRPO to the GSM8K
\cite{cobbe2021training} mathematical reasoning benchmark and to the
Salesforce xLAM-Function-Calling-60k \cite{liu2024xlam} tool-use corpus,
(ii) a set of LoRA adapters published on the HuggingFace Hub
(\texttt{arvindcr4/tinker-rl-bench-*}), and (iii) diagnostic code for
Zero-Variance Fraction, reward-stability, and length-bias analysis. We
analyse four concrete misuse pathways these artefacts could, in principle,
enable, and we describe the existing guardrails and the residual risk we
judge acceptable for publication.

\paragraph{M1. Weaponisation of mathematical reasoning.}
The most capability-relevant artefact we release is GRPO training code evaluated
on grade-school arithmetic and multi-step reasoning; the reviewed evidence is
inconclusive about a capability improvement. A malicious
operator could attempt to extend this pipeline to tasks of greater
dual-use concern, e.g., chemistry olympiad problems, cryptographic key
recovery, or financial fraud. We judge the \emph{marginal} uplift provided
by our release to be small for three reasons. First, GRPO at our training
scale does not create new capabilities; the pre-RL held-out accuracy of
Qwen3-8B-Instruct for GSM8K (Section~\ref{sec:results}, 82.0\%) shows that the
bulk of held-out accuracy is attributable to the checkpoint's pretraining and
instruction tuning. Our 5-seed one-sample seed-level comparison against a fixed
baseline estimates \(+1.3\) percentage points ($p{=}0.26$); it is neither
paired-seed nor reviewed-record evidence of improvement, and should not be read
as a bound on capability change: the capability difference is unmeasured within
wide uncertainty (single base pass, no baseline variance), so ``inconclusive''
does not establish ``small uplift.''
Second, the public literature already contains GRPO implementations
\cite{shao2024deepseekmath, openrlhf2024, verl2024, vonwerra2022trl} that are at
least as capable. Third, any reasoning uplift is bounded by the rewarder:
GSM8K rewards use exact-match against a human-written gold answer, which
does not generalise to the domains listed above without a new reward
signal, which itself requires expertise and labelled data.

\paragraph{M2. Reward hacking and the long-term alignment risk.}
GRPO, like all policy-gradient RL from a proxy reward, is vulnerable to
reward hacking: the model learns to exploit idiosyncrasies of the reward
function rather than solve the intended task \cite{skalse2022specificationgaming,
krakovna2020specification}. Our length-bias discussion
(Section~\ref{sec:length_bias}) records suspicious reward trajectories, not
a verified causal demonstration of reward hacking.
Users who apply our scripts to under-specified reward functions in
production settings (e.g., a custom ``helpfulness'' classifier) may observe
models that appear to improve on held-out metrics while silently
optimising against unintended proxies. We include the Zero-Variance
Fraction and reward-stability diagnostics precisely as retrospective summaries; they do not identify policy drift or reward hacking.

\paragraph{M3. Circumventing safety training via distillation or tool-use.}
The tool-use track trains LoRA adapters that teach a base model to emit
schema-valid JSON function calls \cite{liu2024xlam}. A motivated adversary
could in principle use the same pipeline to teach a base model to emit
jailbreak prompts as ``tool calls'' or to invoke an adversarial external
tool. We mitigate this by (a) releasing only adapters, not merged weights,
while noting that an adapter can still alter safety behavior and does not
guarantee that instruct-model safeguards remain effective; (b) documenting in model
cards that the tool-use adapters are trained on a narrow schema (five
synthetic tools) and have not been safety-evaluated for open-ended function
calling; and (c) not releasing any function-calling harness that would
actually execute emitted calls.

\paragraph{M4. Compute-efficiency externalities.}
This audit makes no demonstrated efficiency claim or group-size recommendation. Improved training efficiency is
itself dual-use: it lowers the cost of running adversarial fine-tuning at
scale. We judge this externality to be modest relative to the already-low
cost of fine-tuning a 0.6B--8B model with commodity LoRA recipes
\cite{hu2022lora}, but we flag it explicitly.

\paragraph{Residual risk.}
After weighing M1--M4 we judge that the reproducibility, calibration, and
pedagogical benefits of releasing \ours{} outweigh the residual misuse
risk. We adopt the standard mitigation of releasing adapters (not merged
weights), documenting intended use and out-of-scope use in model cards,
and publishing alongside the diagnostics a researcher would need to detect
reward hacking.

% ---------------------------------------------------------------------------
\subsection{Compute Cost Accounting}
\label{sec:compute_accounting}
% ---------------------------------------------------------------------------

Table~\ref{tab:compute_spend} retains historical platform entries for audit,
not verified project accounting. Excluding Colab, their numeric ranges sum to
USD 147--177, inconsistent with the earlier USD 130--140 total. Colab is an
additional per-person charge with an unspecified count. The main-text subset
also previously differed. Cash versus credits, dates, cohort overlap, and
failed-run inclusion have not been reconciled; no project spend total is certified.

\begin{table}[ht]
\centering
\small
\caption{Historical compute-spend entries; unreconciled cohorts and cash/credit status.
The sum is an arithmetic check, not verified expenditure \cite{henderson2018deep}.}
\label{tab:compute_spend}
\begin{tabular}{llrr}
\toprule
\textbf{Platform} & \textbf{Use} & \textbf{Runs} & \textbf{Cost (USD)} \\
\midrule
Tinker SDK v0.16.1 & GRPO on 4B/8B (original suite)        & 25 & \$40--55 \\
Tinker SDK v0.16.1 & 10x Structural Ceiling ablation       & 32 & \$65 \\
Tinker SDK v0.16.1 & World-Class Suite (frontier/MoE)      & 13 & \$25--30 \\
Modal H100         & PPO baselines (Qwen3-8B, Llama-3.1-8B)& 2  & \$12--18 \\
Modal H100         & KL-tracking run (failed)              & 1  & \$2--4   \\
Google Colab Pro (T4) & 0.5B--3B QLoRA SFT+GRPO            & --- & \$10/person \\
NVIDIA L4 (TRL baseline)& Qwen2.5-0.5B GSM8K, 5 seeds      & 5  & \$3--5   \\
HuggingFace Hub    & Model + adapter hosting               & --- & \$0 \\
Weights \& Biases  & Experiment tracking (academic tier)   & --- & \$0 \\
\midrule
\multicolumn{3}{l}{\textbf{Project total (unreconciled)}} & \textbf{not available} \\
\bottomrule
\end{tabular}
\end{table}

\paragraph{Hidden costs.} Beyond direct platform spend, the project
consumed human labour (approximately six student-FTE-months, unpaid) and
infrastructure generously subsidised by PES University's LLM Research
Group (shared access to one A100 node used for TRL baselines and
pre-submission sanity checks). No author received commercial compensation
from Thinking Machines, Modal, or any other vendor mentioned.

\paragraph{Failed-run accounting.} Of the Tinker spend, we estimate
\textasciitilde\$30--35 was consumed by runs that ultimately failed
(JWT expiry, API stalls for GPT-OSS-20b and Kimi-K2 before re-run,
KL-tracking gradient bug). We disclose this because reviewers often ask
whether reported total spend reflects only clean, successful runs; it does
not. The failed-run subset is not sufficiently reconciled to estimate a verified
percentage of total spend.

% ---------------------------------------------------------------------------
\subsection{Carbon Footprint Estimation}
\label{sec:carbon}
% ---------------------------------------------------------------------------

Table~\ref{tab:carbon2} preserves conditional energy/carbon scenarios
from historical assumptions, not measured emissions or a project total. We follow the methodology of
\citet{patterson2024empirical} and \citet{strubell2019energy}: for each
GPU-hour we multiply device TDP by wall-clock GPU-hours and the data-centre
Power Usage Effectiveness (PUE), then multiply by the grid carbon intensity
in the region where the hardware is believed to be located.

\paragraph{Assumed constants.}
\begin{itemize}
  \item NVIDIA H100 SXM5 TDP: 700\,W.
  \item NVIDIA A100 80GB TDP: 400\,W.
  \item NVIDIA L4 TDP: 72\,W.
  \item NVIDIA T4 TDP: 70\,W.
  \item Data-centre PUE: 1.1 (conservative; typical hyperscale PUE is
        1.10--1.15 \cite{google2023pue, aws2023pue}).
  \item US average grid intensity (EPA eGRID 2023): 0.367\,kg\,CO$_2$-eq / kWh
        \cite{epa2023egrid}.
  \item Indian average grid intensity (CEA 2023): 0.716\,kg\,CO$_2$-eq / kWh
        \cite{cea2023co2}.
  \item We use the US average for Modal H100 (US-East region as configured)
        and Tinker (location undisclosed; assumed US), and the Indian
        average for Colab Pro T4 used by authors based in Bengaluru
        (Google Cloud asia-south1).
\end{itemize}

\paragraph{GPU-hour estimates.}
Tinker GPU-hour counts are not disclosed by the platform. We estimate them
from wall-clock run duration and assume a single H100-class accelerator per
job, which is an unverified hardware-allocation assumption. The 18 Modal GPU-hours
in this scenario conflict with the 250 GPU-hours in Appendix~\ref{app:compute};
neither is selected as correct without a cohort-specific billing ledger.

\begin{table}[ht]
\centering
\small
\caption{Historical conditional carbon scenarios, not measured emissions. Tinker hardware
is undisclosed; 1$\times$ H100 per run is an assumption. Cohort boundaries and
failed-run coverage are unresolved; the scenario total is not a project estimate.}
\label{tab:carbon2}
\begin{tabular}{lrrrrr}
\toprule
\textbf{Platform} & \textbf{GPU-h} & \textbf{TDP\,(W)} & \textbf{PUE} &
\textbf{Energy\,(kWh)} & \textbf{CO$_2$\,(kg)} \\
\midrule
Tinker (assumed H100)   & \textasciitilde 950 & 700 & 1.1 & \textasciitilde 732 & \textasciitilde 269 \\
Modal H100 (US-East)    & \textasciitilde 18  & 700 & 1.1 & \textasciitilde 14  & \textasciitilde 5.1 \\
PES A100 (in-kind)      & \textasciitilde 60  & 400 & 1.1 & \textasciitilde 26  & \textasciitilde 19   \\
Colab Pro T4 (asia-south1) & \textasciitilde 40 & 70 & 1.2 & \textasciitilde 3.4 & \textasciitilde 2.4  \\
NVIDIA L4 (TRL baseline)   & \textasciitilde 10 & 72 & 1.1 & \textasciitilde 0.8 & \textasciitilde 0.3  \\
\midrule
\textbf{Scenario sum (not project total)} & & & & \textasciitilde 776 & \textasciitilde 296 \\
\bottomrule
\end{tabular}
\end{table}

\paragraph{Interpretation and uncertainty.}
The tabulated sum is conditional on unreconciled hours, assumed accelerator
counts, TDP, PUE, and locations. It is not a verified footprint or an upper bound:
unknown allocation can raise or lower actual consumption. Omitting failed jobs
in a future replication does not establish a bound on this project. A release
needs a single ledger with dates, cohort IDs, GPU counts, measured versus assumed
hours, credits versus cash, and failure status before reporting totals.

\paragraph{Offset and mitigation.}
We did not purchase carbon offsets. Instead, we have (a) made available selected
checkpoint artifacts so that downstream users do not
need to re-run our experiments; (b) retained local step-level CSV logs so that
learning curves can be inspected without re-training; and (c) documented
in Section~\ref{sec:compute_accounting} which runs were wasteful, so that
replicators can skip them. We regard reproducibility-by-artefact as a
more durable mitigation than offsets.

% ---------------------------------------------------------------------------
\subsection{Data Provenance}
\label{sec:provenance}
% ---------------------------------------------------------------------------

\ours{} uses only publicly released research datasets. No private data,
personally identifiable information, or licensed proprietary content is
used. No human annotators were employed during this work. We document
each dataset below.

\paragraph{GSM8K \cite{cobbe2021training}.}
Grade-school math word problems (7{,}473 train / 1{,}319 test)
authored by human writers contracted by OpenAI. Released under the
\textbf{MIT License} via \url{https://github.com/openai/grade-school-math}.
We use the standard train/test splits without modification. The canonical
citation is \citet{cobbe2021training}. Known limitations of GSM8K include
gender-neutral but culturally US-centric word problems (names, currencies,
sports) and a moderate rate of human-labelling errors estimated at
\textasciitilde 2\% by \cite{zhang2024gsm1k}. Our held-out evaluation
respects the test/train split.

\paragraph{Salesforce xLAM-Function-Calling-60k \cite{liu2024xlam}.}
60{,}000 function-calling examples generated by Salesforce Research as
training data for the xLAM agentic model family. Released under the
\textbf{CC BY 4.0} license via \url{https://huggingface.co/datasets/Salesforce/xlam-function-calling-60k}.
We use the public release as-is; specifically, we use the first 35 prompts
$\times$ 10 rollouts in the 10x Structural Ceiling tool-use track and the
full 60k split for the xlam-60k real-data run in Section~\ref{sec:results}.
Known limitations: xLAM schemas are synthetic and skew toward cleanly-typed
arguments; the dataset under-represents ambiguous, error-handling, and
multi-turn tool-calling. Attribution to Salesforce is required by the
license and is provided in Section~\ref{sec:acknowledgments} and in the
xLAM-derived model cards.

\paragraph{HumanEval \cite{chen2021codex}.}
164 Python programming problems with unit tests, released by OpenAI under
the \textbf{MIT License}. The local restricted sandbox does not implement an unmodified standard
pass@$k$ evaluation; code results are harness-bound reward probes. Known
limitations: small scale, English-language docstrings, single-language
coverage; documented contamination concerns against frontier pretraining
corpora \cite{riddell2024contamination}.

\paragraph{NuminaMath \cite{li2024numinamath}.}
Used by collaborator Mohammad Rafi for the multi-stage
GSM8K+NuminaMath pipeline. Released under \textbf{Apache 2.0}. We use a
downstream checkpoint reported by Rafi; our repository contains only his
scripts and the Apache-licensed derivative weights.

\paragraph{Open-Platypus \cite{lee2023platypus}.}
3{,}000 SFT examples used by collaborator Madhu for Qwen3-8B code
generation warm-up. Open-Platypus is a curated subset released under
\textbf{CC BY-NC 4.0}; the non-commercial restriction is respected in our
release, which is academic-use only.

\paragraph{Synthetic tool-use corpus (authored).}
We additionally generated a small five-tool synthetic corpus (calculator,
web-search stub, calendar, file-read, email-send) used for the reviewed descriptive
all-wrong observations. The corpus is authored by the paper's authors from
publicly documented APIs, contains no third-party content, and is
released under the \textbf{MIT License} in our repository under
\texttt{data/synthetic\_tools/}. Generation prompts are included for
auditability. No user data, scraped web content, or proprietary API
traces are present.

\paragraph{No web scraping, no human subjects.}
We did not scrape any website. We did not run any study with human
participants. No IRB review was therefore required. No data that could
reasonably be considered to contain PII, copyrighted content, or sensitive
personal attributes was used at any stage of training, evaluation, or
reward modelling.

% ---------------------------------------------------------------------------
\subsection{Closed-Source Tinker Acknowledgment}
\label{sec:tinker_closed}
% ---------------------------------------------------------------------------

A central tension in \ours{} is that a substantial fraction of our
headline numbers come from a closed-source commercial platform while our
paper simultaneously advocates for reproducibility and platform
independence. We do not resolve this tension by hiding it; we describe
it precisely.

\paragraph{What Tinker is and is not.}
Tinker \cite{thinkingmachines2024tinker} is a managed LLM fine-tuning and
inference service provided by Thinking Machines, Inc. It exposes a Python
SDK (\texttt{tinker==0.16.1}) that accepts custom loss functions
(\texttt{forward\_backward\_custom}), a limited set of optimisers, and
standard LoRA hyperparameters. Local scripts specify reward normalization, baseline subtraction, custom loss,
and requested LoRA/optimizer configuration. Undisclosed server-side details
include hardware placement, scheduling, execution numerics, and some batching
and sampler behavior. These are separate from the inspectable local objective.

\paragraph{What this means for our claims.}
We have not ruled out seed variance, checkpoint/model-size confounding, evaluator
differences, or backend effects. The arithmetic/GSM8K cross-library identity is
unresolved. Tinker results do not establish an upper bound or algorithmic
superiority. verl/OpenRLHF Task~4 dry runs are engineering artifacts, not
measured alternatives. The contribution remains the centered-reward identity,
descriptive audited cells, and reporting contract.

\paragraph{Reproducibility commitments we can make.}
Section~\ref{sec:reproducibility} distinguishes local source, observed
configuration, undisclosed internals, and evidence access. A dagger denotes only
interrupted execution; backend provenance is stated explicitly. Exact checkpoint
revisions and public access must be verified per artifact; no blanket release
or future rerun commitment is made here.

\paragraph{Key rotation and credential hygiene.}
Our Tinker API key (prefix \texttt{tml-...}) is held in a repository
\texttt{.env.example} placeholder only. The real key is stored in the
authors' password manager and rotated whenever a failure mode suggests
possible exfiltration. No Tinker key is committed to git history in
either the primary repository (\texttt{pes-llm-research/tinker-rl-lab}) or
the mirror (\texttt{arvindcr4/tinker-rl-lab}).

% ---------------------------------------------------------------------------
\subsection{Known Methodological Limits}
\label{sec:methodlimits}
% ---------------------------------------------------------------------------

In addition to the infrastructure failures and platform-specific caveats
enumerated in Section~\ref{sec:limitations}, we flag the following
methodological limits.

\paragraph{Short training horizons.}
Most reviewed Tinker runs have short, often 30-step horizons; this is not
a universal limit for all later scripts or campaigns. This is a
deliberate cost-control choice and means our Tinker numbers are
\emph{early-training snapshots}. We cannot rule out that longer training
would change the qualitative picture. The later 82\%$\to$83.3\% comparison
is one-sample fixed-baseline evidence, not paired-seed or reviewed-record
evidence; it is not a held-out gain whose size can be extrapolated.

\paragraph{Single-seed Tinker experiments.}
Each Tinker configuration was run once, not 3--5 times as best practice
prescribes \cite{henderson2018deep}. Tinker results therefore lack
variance estimates and do not support significance testing; we report
them descriptively. Seed-level inference requires the actual independent seed observations and
an identified estimand; reconstructed seed vectors are not eligible.

\paragraph{Train-set reward as primary Tinker metric.}
Tinker runs primarily report reward on training prompts. Only the GSM8K
track was followed up with a separate 200-example held-out evaluation
per seed. Other Tinker tracks (tool-use, xLAM) remain train-set-only and
we cannot distinguish memorisation from generalisation for those
settings.

\paragraph{LoRA only; no full fine-tuning.}
The reviewed LLM fine-tuning runs use LoRA \cite{hu2022lora}, with settings
varying by cohort; the Wave~6 rank sweep also includes rank 4. We have
\emph{not} tested whether full fine-tuning would re-order the
library/algorithm comparisons. The LoRA constraint is practical (cost) but
it means our claims about PPO vs.\ GRPO and TRL vs.\ Tinker are
LoRA-conditional.

\paragraph{Benchmark coverage is narrow.}
We evaluate four domains: GSM8K, MATH-500 (exploratory), HumanEval
(subset), and synthetic/xLAM tool-use. We do not evaluate on MT-Bench,
ArenaHard, safety benchmarks (HarmBench, ToxicChat), or truthfulness
benchmarks (TruthfulQA). Any claim about ``reasoning'' in the paper is
strictly a claim about the covered benchmarks.

\paragraph{No human preference data.}
We use verifiable rewards only (exact-match for math, unit-test for code,
schema-validity for tool-use). We do not study RLHF with human preference
data; findings should not be extrapolated to reward-model--based alignment
without further work.

\paragraph{Closed-source implementation opacity (reiterated).}
Comparisons between Tinker GRPO and TRL GRPO are confounded by the
closed-source nature of Tinker's implementation. See
Section~\ref{sec:tinker_closed} for detail.

\paragraph{Cross-platform hardware confounding.}
Tinker, Modal, Colab, and the PES A100 node use different accelerators,
different memory hierarchies, and different inference/training stacks
(Tinker proprietary, Modal vLLM, Colab PyTorch-native, TRL+Accelerate).
Observed differences in reward dynamics are not purely algorithmic.

\paragraph{Carbon estimates are approximate.}
As discussed in Section~\ref{sec:carbon}, Tinker GPU-hour and TDP are
inferred, not measured. Grid intensity is region-average, not marginal.
Numbers in Table~\ref{tab:carbon2} are conditional scenarios with
unreconciled cohorts, not measured emissions or upper bounds.

\paragraph{Exploratory, not confirmatory.}
This is an exploratory case study, not a pre-registered confirmatory
experiment. We did not pre-register primary hypotheses. The conflicting historical correction families and step-based significance
decisions are withdrawn (Section~\ref{sec:statistical-policy}); exploratory
coefficients do not establish calibrated inference.

\paragraph{Geographic and demographic limits.}
All authors are based at a single institution in Bengaluru, India. Our
choice of benchmarks, prompt phrasings, and evaluation design reflects
that context. Independent replication by groups in other regions, on
non-English tasks, and with non-Qwen / non-Llama families is an open
question.

% =============================================================================
% End of ethics_statement.tex
 into main.tex
% Intended placement: after the Conclusion section, before References.
% Label: \label{sec:impact_standalone}
% Target length: >= 2 pages (NeurIPS Broader Impact guidance).
% =============================================================================

\section{Ethics, Limitations, and Broader Impact}
\label{sec:impact_standalone}

This statement is written to satisfy the NeurIPS Code of Ethics and Broader
Impact requirements. It complements the shorter Limitations and Broader
Impact subsections embedded in Section~\ref{sec:limitations} by consolidating
in one place a dual-use analysis, unresolved compute-accounting discrepancies, carbon
scenarios, data provenance disclosures, a candid acknowledgment
of the closed-source training backend on which a fraction of our headline
numbers depends, and a list of methodological limits we are aware of but
did not have resources to close within the submission window.

% ---------------------------------------------------------------------------
\subsection{Dual-Use Analysis: Misuse Risks of Reasoning RL}
\label{sec:dualuse}
% ---------------------------------------------------------------------------

\paragraph{Threat model.}
\ours{} releases (i) training scripts that apply GRPO to the GSM8K
\cite{cobbe2021training} mathematical reasoning benchmark and to the
Salesforce xLAM-Function-Calling-60k \cite{liu2024xlam} tool-use corpus,
(ii) a set of LoRA adapters published on the HuggingFace Hub
(\texttt{arvindcr4/tinker-rl-bench-*}), and (iii) diagnostic code for
Zero-Variance Fraction, reward-stability, and length-bias analysis. We
analyse four concrete misuse pathways these artefacts could, in principle,
enable, and we describe the existing guardrails and the residual risk we
judge acceptable for publication.

\paragraph{M1. Weaponisation of mathematical reasoning.}
The most capability-relevant artefact we release is GRPO training code evaluated
on grade-school arithmetic and multi-step reasoning; the reviewed evidence is
inconclusive about a capability improvement. A malicious
operator could attempt to extend this pipeline to tasks of greater
dual-use concern, e.g., chemistry olympiad problems, cryptographic key
recovery, or financial fraud. We judge the \emph{marginal} uplift provided
by our release to be small for three reasons. First, GRPO at our training
scale does not create new capabilities; the pre-RL held-out accuracy of
Qwen3-8B-Instruct for GSM8K (Section~\ref{sec:results}, 82.0\%) shows that the
bulk of held-out accuracy is attributable to the checkpoint's pretraining and
instruction tuning. Our 5-seed one-sample seed-level comparison against a fixed
baseline estimates \(+1.3\) percentage points ($p{=}0.26$); it is neither
paired-seed nor reviewed-record evidence of improvement, and should not be read
as a bound on capability change: the capability difference is unmeasured within
wide uncertainty (single base pass, no baseline variance), so ``inconclusive''
does not establish ``small uplift.''
Second, the public literature already contains GRPO implementations
\cite{shao2024deepseekmath, openrlhf2024, verl2024, vonwerra2022trl} that are at
least as capable. Third, any reasoning uplift is bounded by the rewarder:
GSM8K rewards use exact-match against a human-written gold answer, which
does not generalise to the domains listed above without a new reward
signal, which itself requires expertise and labelled data.

\paragraph{M2. Reward hacking and the long-term alignment risk.}
GRPO, like all policy-gradient RL from a proxy reward, is vulnerable to
reward hacking: the model learns to exploit idiosyncrasies of the reward
function rather than solve the intended task \cite{skalse2022specificationgaming,
krakovna2020specification}. Our length-bias discussion
(Section~\ref{sec:length_bias}) records suspicious reward trajectories, not
a verified causal demonstration of reward hacking.
Users who apply our scripts to under-specified reward functions in
production settings (e.g., a custom ``helpfulness'' classifier) may observe
models that appear to improve on held-out metrics while silently
optimising against unintended proxies. We include the Zero-Variance
Fraction and reward-stability diagnostics precisely as retrospective summaries; they do not identify policy drift or reward hacking.

\paragraph{M3. Circumventing safety training via distillation or tool-use.}
The tool-use track trains LoRA adapters that teach a base model to emit
schema-valid JSON function calls \cite{liu2024xlam}. A motivated adversary
could in principle use the same pipeline to teach a base model to emit
jailbreak prompts as ``tool calls'' or to invoke an adversarial external
tool. We mitigate this by (a) releasing only adapters, not merged weights,
while noting that an adapter can still alter safety behavior and does not
guarantee that instruct-model safeguards remain effective; (b) documenting in model
cards that the tool-use adapters are trained on a narrow schema (five
synthetic tools) and have not been safety-evaluated for open-ended function
calling; and (c) not releasing any function-calling harness that would
actually execute emitted calls.

\paragraph{M4. Compute-efficiency externalities.}
This audit makes no demonstrated efficiency claim or group-size recommendation. Improved training efficiency is
itself dual-use: it lowers the cost of running adversarial fine-tuning at
scale. We judge this externality to be modest relative to the already-low
cost of fine-tuning a 0.6B--8B model with commodity LoRA recipes
\cite{hu2022lora}, but we flag it explicitly.

\paragraph{Residual risk.}
After weighing M1--M4 we judge that the reproducibility, calibration, and
pedagogical benefits of releasing \ours{} outweigh the residual misuse
risk. We adopt the standard mitigation of releasing adapters (not merged
weights), documenting intended use and out-of-scope use in model cards,
and publishing alongside the diagnostics a researcher would need to detect
reward hacking.

% ---------------------------------------------------------------------------
\subsection{Compute Cost Accounting}
\label{sec:compute_accounting}
% ---------------------------------------------------------------------------

Table~\ref{tab:compute_spend} retains historical platform entries for audit,
not verified project accounting. Excluding Colab, their numeric ranges sum to
USD 147--177, inconsistent with the earlier USD 130--140 total. Colab is an
additional per-person charge with an unspecified count. The main-text subset
also previously differed. Cash versus credits, dates, cohort overlap, and
failed-run inclusion have not been reconciled; no project spend total is certified.

\begin{table}[ht]
\centering
\small
\caption{Historical compute-spend entries; unreconciled cohorts and cash/credit status.
The sum is an arithmetic check, not verified expenditure \cite{henderson2018deep}.}
\label{tab:compute_spend}
\begin{tabular}{llrr}
\toprule
\textbf{Platform} & \textbf{Use} & \textbf{Runs} & \textbf{Cost (USD)} \\
\midrule
Tinker SDK v0.16.1 & GRPO on 4B/8B (original suite)        & 25 & \$40--55 \\
Tinker SDK v0.16.1 & 10x Structural Ceiling ablation       & 32 & \$65 \\
Tinker SDK v0.16.1 & World-Class Suite (frontier/MoE)      & 13 & \$25--30 \\
Modal H100         & PPO baselines (Qwen3-8B, Llama-3.1-8B)& 2  & \$12--18 \\
Modal H100         & KL-tracking run (failed)              & 1  & \$2--4   \\
Google Colab Pro (T4) & 0.5B--3B QLoRA SFT+GRPO            & --- & \$10/person \\
NVIDIA L4 (TRL baseline)& Qwen2.5-0.5B GSM8K, 5 seeds      & 5  & \$3--5   \\
HuggingFace Hub    & Model + adapter hosting               & --- & \$0 \\
Weights \& Biases  & Experiment tracking (academic tier)   & --- & \$0 \\
\midrule
\multicolumn{3}{l}{\textbf{Project total (unreconciled)}} & \textbf{not available} \\
\bottomrule
\end{tabular}
\end{table}

\paragraph{Hidden costs.} Beyond direct platform spend, the project
consumed human labour (approximately six student-FTE-months, unpaid) and
infrastructure generously subsidised by PES University's LLM Research
Group (shared access to one A100 node used for TRL baselines and
pre-submission sanity checks). No author received commercial compensation
from Thinking Machines, Modal, or any other vendor mentioned.

\paragraph{Failed-run accounting.} Of the Tinker spend, we estimate
\textasciitilde\$30--35 was consumed by runs that ultimately failed
(JWT expiry, API stalls for GPT-OSS-20b and Kimi-K2 before re-run,
KL-tracking gradient bug). We disclose this because reviewers often ask
whether reported total spend reflects only clean, successful runs; it does
not. The failed-run subset is not sufficiently reconciled to estimate a verified
percentage of total spend.

% ---------------------------------------------------------------------------
\subsection{Carbon Footprint Estimation}
\label{sec:carbon}
% ---------------------------------------------------------------------------

Table~\ref{tab:carbon2} preserves conditional energy/carbon scenarios
from historical assumptions, not measured emissions or a project total. We follow the methodology of
\citet{patterson2024empirical} and \citet{strubell2019energy}: for each
GPU-hour we multiply device TDP by wall-clock GPU-hours and the data-centre
Power Usage Effectiveness (PUE), then multiply by the grid carbon intensity
in the region where the hardware is believed to be located.

\paragraph{Assumed constants.}
\begin{itemize}
  \item NVIDIA H100 SXM5 TDP: 700\,W.
  \item NVIDIA A100 80GB TDP: 400\,W.
  \item NVIDIA L4 TDP: 72\,W.
  \item NVIDIA T4 TDP: 70\,W.
  \item Data-centre PUE: 1.1 (conservative; typical hyperscale PUE is
        1.10--1.15 \cite{google2023pue, aws2023pue}).
  \item US average grid intensity (EPA eGRID 2023): 0.367\,kg\,CO$_2$-eq / kWh
        \cite{epa2023egrid}.
  \item Indian average grid intensity (CEA 2023): 0.716\,kg\,CO$_2$-eq / kWh
        \cite{cea2023co2}.
  \item We use the US average for Modal H100 (US-East region as configured)
        and Tinker (location undisclosed; assumed US), and the Indian
        average for Colab Pro T4 used by authors based in Bengaluru
        (Google Cloud asia-south1).
\end{itemize}

\paragraph{GPU-hour estimates.}
Tinker GPU-hour counts are not disclosed by the platform. We estimate them
from wall-clock run duration and assume a single H100-class accelerator per
job, which is an unverified hardware-allocation assumption. The 18 Modal GPU-hours
in this scenario conflict with the 250 GPU-hours in Appendix~\ref{app:compute};
neither is selected as correct without a cohort-specific billing ledger.

\begin{table}[ht]
\centering
\small
\caption{Historical conditional carbon scenarios, not measured emissions. Tinker hardware
is undisclosed; 1$\times$ H100 per run is an assumption. Cohort boundaries and
failed-run coverage are unresolved; the scenario total is not a project estimate.}
\label{tab:carbon2}
\begin{tabular}{lrrrrr}
\toprule
\textbf{Platform} & \textbf{GPU-h} & \textbf{TDP\,(W)} & \textbf{PUE} &
\textbf{Energy\,(kWh)} & \textbf{CO$_2$\,(kg)} \\
\midrule
Tinker (assumed H100)   & \textasciitilde 950 & 700 & 1.1 & \textasciitilde 732 & \textasciitilde 269 \\
Modal H100 (US-East)    & \textasciitilde 18  & 700 & 1.1 & \textasciitilde 14  & \textasciitilde 5.1 \\
PES A100 (in-kind)      & \textasciitilde 60  & 400 & 1.1 & \textasciitilde 26  & \textasciitilde 19   \\
Colab Pro T4 (asia-south1) & \textasciitilde 40 & 70 & 1.2 & \textasciitilde 3.4 & \textasciitilde 2.4  \\
NVIDIA L4 (TRL baseline)   & \textasciitilde 10 & 72 & 1.1 & \textasciitilde 0.8 & \textasciitilde 0.3  \\
\midrule
\textbf{Scenario sum (not project total)} & & & & \textasciitilde 776 & \textasciitilde 296 \\
\bottomrule
\end{tabular}
\end{table}

\paragraph{Interpretation and uncertainty.}
The tabulated sum is conditional on unreconciled hours, assumed accelerator
counts, TDP, PUE, and locations. It is not a verified footprint or an upper bound:
unknown allocation can raise or lower actual consumption. Omitting failed jobs
in a future replication does not establish a bound on this project. A release
needs a single ledger with dates, cohort IDs, GPU counts, measured versus assumed
hours, credits versus cash, and failure status before reporting totals.

\paragraph{Offset and mitigation.}
We did not purchase carbon offsets. Instead, we have (a) made available selected
checkpoint artifacts so that downstream users do not
need to re-run our experiments; (b) retained local step-level CSV logs so that
learning curves can be inspected without re-training; and (c) documented
in Section~\ref{sec:compute_accounting} which runs were wasteful, so that
replicators can skip them. We regard reproducibility-by-artefact as a
more durable mitigation than offsets.

% ---------------------------------------------------------------------------
\subsection{Data Provenance}
\label{sec:provenance}
% ---------------------------------------------------------------------------

\ours{} uses only publicly released research datasets. No private data,
personally identifiable information, or licensed proprietary content is
used. No human annotators were employed during this work. We document
each dataset below.

\paragraph{GSM8K \cite{cobbe2021training}.}
Grade-school math word problems (7{,}473 train / 1{,}319 test)
authored by human writers contracted by OpenAI. Released under the
\textbf{MIT License} via \url{https://github.com/openai/grade-school-math}.
We use the standard train/test splits without modification. The canonical
citation is \citet{cobbe2021training}. Known limitations of GSM8K include
gender-neutral but culturally US-centric word problems (names, currencies,
sports) and a moderate rate of human-labelling errors estimated at
\textasciitilde 2\% by \cite{zhang2024gsm1k}. Our held-out evaluation
respects the test/train split.

\paragraph{Salesforce xLAM-Function-Calling-60k \cite{liu2024xlam}.}
60{,}000 function-calling examples generated by Salesforce Research as
training data for the xLAM agentic model family. Released under the
\textbf{CC BY 4.0} license via \url{https://huggingface.co/datasets/Salesforce/xlam-function-calling-60k}.
We use the public release as-is; specifically, we use the first 35 prompts
$\times$ 10 rollouts in the 10x Structural Ceiling tool-use track and the
full 60k split for the xlam-60k real-data run in Section~\ref{sec:results}.
Known limitations: xLAM schemas are synthetic and skew toward cleanly-typed
arguments; the dataset under-represents ambiguous, error-handling, and
multi-turn tool-calling. Attribution to Salesforce is required by the
license and is provided in Section~\ref{sec:acknowledgments} and in the
xLAM-derived model cards.

\paragraph{HumanEval \cite{chen2021codex}.}
164 Python programming problems with unit tests, released by OpenAI under
the \textbf{MIT License}. The local restricted sandbox does not implement an unmodified standard
pass@$k$ evaluation; code results are harness-bound reward probes. Known
limitations: small scale, English-language docstrings, single-language
coverage; documented contamination concerns against frontier pretraining
corpora \cite{riddell2024contamination}.

\paragraph{NuminaMath \cite{li2024numinamath}.}
Used by collaborator Mohammad Rafi for the multi-stage
GSM8K+NuminaMath pipeline. Released under \textbf{Apache 2.0}. We use a
downstream checkpoint reported by Rafi; our repository contains only his
scripts and the Apache-licensed derivative weights.

\paragraph{Open-Platypus \cite{lee2023platypus}.}
3{,}000 SFT examples used by collaborator Madhu for Qwen3-8B code
generation warm-up. Open-Platypus is a curated subset released under
\textbf{CC BY-NC 4.0}; the non-commercial restriction is respected in our
release, which is academic-use only.

\paragraph{Synthetic tool-use corpus (authored).}
We additionally generated a small five-tool synthetic corpus (calculator,
web-search stub, calendar, file-read, email-send) used for the reviewed descriptive
all-wrong observations. The corpus is authored by the paper's authors from
publicly documented APIs, contains no third-party content, and is
released under the \textbf{MIT License} in our repository under
\texttt{data/synthetic\_tools/}. Generation prompts are included for
auditability. No user data, scraped web content, or proprietary API
traces are present.

\paragraph{No web scraping, no human subjects.}
We did not scrape any website. We did not run any study with human
participants. No IRB review was therefore required. No data that could
reasonably be considered to contain PII, copyrighted content, or sensitive
personal attributes was used at any stage of training, evaluation, or
reward modelling.

% ---------------------------------------------------------------------------
\subsection{Closed-Source Tinker Acknowledgment}
\label{sec:tinker_closed}
% ---------------------------------------------------------------------------

A central tension in \ours{} is that a substantial fraction of our
headline numbers come from a closed-source commercial platform while our
paper simultaneously advocates for reproducibility and platform
independence. We do not resolve this tension by hiding it; we describe
it precisely.

\paragraph{What Tinker is and is not.}
Tinker \cite{thinkingmachines2024tinker} is a managed LLM fine-tuning and
inference service provided by Thinking Machines, Inc. It exposes a Python
SDK (\texttt{tinker==0.16.1}) that accepts custom loss functions
(\texttt{forward\_backward\_custom}), a limited set of optimisers, and
standard LoRA hyperparameters. Local scripts specify reward normalization, baseline subtraction, custom loss,
and requested LoRA/optimizer configuration. Undisclosed server-side details
include hardware placement, scheduling, execution numerics, and some batching
and sampler behavior. These are separate from the inspectable local objective.

\paragraph{What this means for our claims.}
We have not ruled out seed variance, checkpoint/model-size confounding, evaluator
differences, or backend effects. The arithmetic/GSM8K cross-library identity is
unresolved. Tinker results do not establish an upper bound or algorithmic
superiority. verl/OpenRLHF Task~4 dry runs are engineering artifacts, not
measured alternatives. The contribution remains the centered-reward identity,
descriptive audited cells, and reporting contract.

\paragraph{Reproducibility commitments we can make.}
Section~\ref{sec:reproducibility} distinguishes local source, observed
configuration, undisclosed internals, and evidence access. A dagger denotes only
interrupted execution; backend provenance is stated explicitly. Exact checkpoint
revisions and public access must be verified per artifact; no blanket release
or future rerun commitment is made here.

\paragraph{Key rotation and credential hygiene.}
Our Tinker API key (prefix \texttt{tml-...}) is held in a repository
\texttt{.env.example} placeholder only. The real key is stored in the
authors' password manager and rotated whenever a failure mode suggests
possible exfiltration. No Tinker key is committed to git history in
either the primary repository (\texttt{pes-llm-research/tinker-rl-lab}) or
the mirror (\texttt{arvindcr4/tinker-rl-lab}).

% ---------------------------------------------------------------------------
\subsection{Known Methodological Limits}
\label{sec:methodlimits}
% ---------------------------------------------------------------------------

In addition to the infrastructure failures and platform-specific caveats
enumerated in Section~\ref{sec:limitations}, we flag the following
methodological limits.

\paragraph{Short training horizons.}
Most reviewed Tinker runs have short, often 30-step horizons; this is not
a universal limit for all later scripts or campaigns. This is a
deliberate cost-control choice and means our Tinker numbers are
\emph{early-training snapshots}. We cannot rule out that longer training
would change the qualitative picture. The later 82\%$\to$83.3\% comparison
is one-sample fixed-baseline evidence, not paired-seed or reviewed-record
evidence; it is not a held-out gain whose size can be extrapolated.

\paragraph{Single-seed Tinker experiments.}
Each Tinker configuration was run once, not 3--5 times as best practice
prescribes \cite{henderson2018deep}. Tinker results therefore lack
variance estimates and do not support significance testing; we report
them descriptively. Seed-level inference requires the actual independent seed observations and
an identified estimand; reconstructed seed vectors are not eligible.

\paragraph{Train-set reward as primary Tinker metric.}
Tinker runs primarily report reward on training prompts. Only the GSM8K
track was followed up with a separate 200-example held-out evaluation
per seed. Other Tinker tracks (tool-use, xLAM) remain train-set-only and
we cannot distinguish memorisation from generalisation for those
settings.

\paragraph{LoRA only; no full fine-tuning.}
The reviewed LLM fine-tuning runs use LoRA \cite{hu2022lora}, with settings
varying by cohort; the Wave~6 rank sweep also includes rank 4. We have
\emph{not} tested whether full fine-tuning would re-order the
library/algorithm comparisons. The LoRA constraint is practical (cost) but
it means our claims about PPO vs.\ GRPO and TRL vs.\ Tinker are
LoRA-conditional.

\paragraph{Benchmark coverage is narrow.}
We evaluate four domains: GSM8K, MATH-500 (exploratory), HumanEval
(subset), and synthetic/xLAM tool-use. We do not evaluate on MT-Bench,
ArenaHard, safety benchmarks (HarmBench, ToxicChat), or truthfulness
benchmarks (TruthfulQA). Any claim about ``reasoning'' in the paper is
strictly a claim about the covered benchmarks.

\paragraph{No human preference data.}
We use verifiable rewards only (exact-match for math, unit-test for code,
schema-validity for tool-use). We do not study RLHF with human preference
data; findings should not be extrapolated to reward-model--based alignment
without further work.

\paragraph{Closed-source implementation opacity (reiterated).}
Comparisons between Tinker GRPO and TRL GRPO are confounded by the
closed-source nature of Tinker's implementation. See
Section~\ref{sec:tinker_closed} for detail.

\paragraph{Cross-platform hardware confounding.}
Tinker, Modal, Colab, and the PES A100 node use different accelerators,
different memory hierarchies, and different inference/training stacks
(Tinker proprietary, Modal vLLM, Colab PyTorch-native, TRL+Accelerate).
Observed differences in reward dynamics are not purely algorithmic.

\paragraph{Carbon estimates are approximate.}
As discussed in Section~\ref{sec:carbon}, Tinker GPU-hour and TDP are
inferred, not measured. Grid intensity is region-average, not marginal.
Numbers in Table~\ref{tab:carbon2} are conditional scenarios with
unreconciled cohorts, not measured emissions or upper bounds.

\paragraph{Exploratory, not confirmatory.}
This is an exploratory case study, not a pre-registered confirmatory
experiment. We did not pre-register primary hypotheses. The conflicting historical correction families and step-based significance
decisions are withdrawn (Section~\ref{sec:statistical-policy}); exploratory
coefficients do not establish calibrated inference.

\paragraph{Geographic and demographic limits.}
All authors are based at a single institution in Bengaluru, India. Our
choice of benchmarks, prompt phrasings, and evaluation design reflects
that context. Independent replication by groups in other regions, on
non-English tasks, and with non-Qwen / non-Llama families is an open
question.

% =============================================================================
% End of ethics_statement.tex
 into main.tex
% Intended placement: after the Conclusion section, before References.
% Label: \label{sec:impact_standalone}
% Target length: >= 2 pages (NeurIPS Broader Impact guidance).
% =============================================================================

\section{Ethics, Limitations, and Broader Impact}
\label{sec:impact_standalone}

This statement is written to satisfy the NeurIPS Code of Ethics and Broader
Impact requirements. It complements the shorter Limitations and Broader
Impact subsections embedded in Section~\ref{sec:limitations} by consolidating
in one place a dual-use analysis, unresolved compute-accounting discrepancies, carbon
scenarios, data provenance disclosures, a candid acknowledgment
of the closed-source training backend on which a fraction of our headline
numbers depends, and a list of methodological limits we are aware of but
did not have resources to close within the submission window.

% ---------------------------------------------------------------------------
\subsection{Dual-Use Analysis: Misuse Risks of Reasoning RL}
\label{sec:dualuse}
% ---------------------------------------------------------------------------

\paragraph{Threat model.}
\ours{} releases (i) training scripts that apply GRPO to the GSM8K
\cite{cobbe2021training} mathematical reasoning benchmark and to the
Salesforce xLAM-Function-Calling-60k \cite{liu2024xlam} tool-use corpus,
(ii) a set of LoRA adapters published on the HuggingFace Hub
(\texttt{arvindcr4/tinker-rl-bench-*}), and (iii) diagnostic code for
Zero-Variance Fraction, reward-stability, and length-bias analysis. We
analyse four concrete misuse pathways these artefacts could, in principle,
enable, and we describe the existing guardrails and the residual risk we
judge acceptable for publication.

\paragraph{M1. Weaponisation of mathematical reasoning.}
The most capability-relevant artefact we release is GRPO training code evaluated
on grade-school arithmetic and multi-step reasoning; the reviewed evidence is
inconclusive about a capability improvement. A malicious
operator could attempt to extend this pipeline to tasks of greater
dual-use concern, e.g., chemistry olympiad problems, cryptographic key
recovery, or financial fraud. We judge the \emph{marginal} uplift provided
by our release to be small for three reasons. First, GRPO at our training
scale does not create new capabilities; the pre-RL held-out accuracy of
Qwen3-8B-Instruct for GSM8K (Section~\ref{sec:results}, 82.0\%) shows that the
bulk of held-out accuracy is attributable to the checkpoint's pretraining and
instruction tuning. Our 5-seed one-sample seed-level comparison against a fixed
baseline estimates \(+1.3\) percentage points ($p{=}0.26$); it is neither
paired-seed nor reviewed-record evidence of improvement, and should not be read
as a bound on capability change: the capability difference is unmeasured within
wide uncertainty (single base pass, no baseline variance), so ``inconclusive''
does not establish ``small uplift.''
Second, the public literature already contains GRPO implementations
\cite{shao2024deepseekmath, openrlhf2024, verl2024, vonwerra2022trl} that are at
least as capable. Third, any reasoning uplift is bounded by the rewarder:
GSM8K rewards use exact-match against a human-written gold answer, which
does not generalise to the domains listed above without a new reward
signal, which itself requires expertise and labelled data.

\paragraph{M2. Reward hacking and the long-term alignment risk.}
GRPO, like all policy-gradient RL from a proxy reward, is vulnerable to
reward hacking: the model learns to exploit idiosyncrasies of the reward
function rather than solve the intended task \cite{skalse2022specificationgaming,
krakovna2020specification}. Our length-bias discussion
(Section~\ref{sec:length_bias}) records suspicious reward trajectories, not
a verified causal demonstration of reward hacking.
Users who apply our scripts to under-specified reward functions in
production settings (e.g., a custom ``helpfulness'' classifier) may observe
models that appear to improve on held-out metrics while silently
optimising against unintended proxies. We include the Zero-Variance
Fraction and reward-stability diagnostics precisely as retrospective summaries; they do not identify policy drift or reward hacking.

\paragraph{M3. Circumventing safety training via distillation or tool-use.}
The tool-use track trains LoRA adapters that teach a base model to emit
schema-valid JSON function calls \cite{liu2024xlam}. A motivated adversary
could in principle use the same pipeline to teach a base model to emit
jailbreak prompts as ``tool calls'' or to invoke an adversarial external
tool. We mitigate this by (a) releasing only adapters, not merged weights,
while noting that an adapter can still alter safety behavior and does not
guarantee that instruct-model safeguards remain effective; (b) documenting in model
cards that the tool-use adapters are trained on a narrow schema (five
synthetic tools) and have not been safety-evaluated for open-ended function
calling; and (c) not releasing any function-calling harness that would
actually execute emitted calls.

\paragraph{M4. Compute-efficiency externalities.}
This audit makes no demonstrated efficiency claim or group-size recommendation. Improved training efficiency is
itself dual-use: it lowers the cost of running adversarial fine-tuning at
scale. We judge this externality to be modest relative to the already-low
cost of fine-tuning a 0.6B--8B model with commodity LoRA recipes
\cite{hu2022lora}, but we flag it explicitly.

\paragraph{Residual risk.}
After weighing M1--M4 we judge that the reproducibility, calibration, and
pedagogical benefits of releasing \ours{} outweigh the residual misuse
risk. We adopt the standard mitigation of releasing adapters (not merged
weights), documenting intended use and out-of-scope use in model cards,
and publishing alongside the diagnostics a researcher would need to detect
reward hacking.

% ---------------------------------------------------------------------------
\subsection{Compute Cost Accounting}
\label{sec:compute_accounting}
% ---------------------------------------------------------------------------

Table~\ref{tab:compute_spend} retains historical platform entries for audit,
not verified project accounting. Excluding Colab, their numeric ranges sum to
USD 147--177, inconsistent with the earlier USD 130--140 total. Colab is an
additional per-person charge with an unspecified count. The main-text subset
also previously differed. Cash versus credits, dates, cohort overlap, and
failed-run inclusion have not been reconciled; no project spend total is certified.

\begin{table}[ht]
\centering
\small
\caption{Historical compute-spend entries; unreconciled cohorts and cash/credit status.
The sum is an arithmetic check, not verified expenditure \cite{henderson2018deep}.}
\label{tab:compute_spend}
\begin{tabular}{llrr}
\toprule
\textbf{Platform} & \textbf{Use} & \textbf{Runs} & \textbf{Cost (USD)} \\
\midrule
Tinker SDK v0.16.1 & GRPO on 4B/8B (original suite)        & 25 & \$40--55 \\
Tinker SDK v0.16.1 & 10x Structural Ceiling ablation       & 32 & \$65 \\
Tinker SDK v0.16.1 & World-Class Suite (frontier/MoE)      & 13 & \$25--30 \\
Modal H100         & PPO baselines (Qwen3-8B, Llama-3.1-8B)& 2  & \$12--18 \\
Modal H100         & KL-tracking run (failed)              & 1  & \$2--4   \\
Google Colab Pro (T4) & 0.5B--3B QLoRA SFT+GRPO            & --- & \$10/person \\
NVIDIA L4 (TRL baseline)& Qwen2.5-0.5B GSM8K, 5 seeds      & 5  & \$3--5   \\
HuggingFace Hub    & Model + adapter hosting               & --- & \$0 \\
Weights \& Biases  & Experiment tracking (academic tier)   & --- & \$0 \\
\midrule
\multicolumn{3}{l}{\textbf{Project total (unreconciled)}} & \textbf{not available} \\
\bottomrule
\end{tabular}
\end{table}

\paragraph{Hidden costs.} Beyond direct platform spend, the project
consumed human labour (approximately six student-FTE-months, unpaid) and
infrastructure generously subsidised by PES University's LLM Research
Group (shared access to one A100 node used for TRL baselines and
pre-submission sanity checks). No author received commercial compensation
from Thinking Machines, Modal, or any other vendor mentioned.

\paragraph{Failed-run accounting.} Of the Tinker spend, we estimate
\textasciitilde\$30--35 was consumed by runs that ultimately failed
(JWT expiry, API stalls for GPT-OSS-20b and Kimi-K2 before re-run,
KL-tracking gradient bug). We disclose this because reviewers often ask
whether reported total spend reflects only clean, successful runs; it does
not. The failed-run subset is not sufficiently reconciled to estimate a verified
percentage of total spend.

% ---------------------------------------------------------------------------
\subsection{Carbon Footprint Estimation}
\label{sec:carbon}
% ---------------------------------------------------------------------------

Table~\ref{tab:carbon2} preserves conditional energy/carbon scenarios
from historical assumptions, not measured emissions or a project total. We follow the methodology of
\citet{patterson2024empirical} and \citet{strubell2019energy}: for each
GPU-hour we multiply device TDP by wall-clock GPU-hours and the data-centre
Power Usage Effectiveness (PUE), then multiply by the grid carbon intensity
in the region where the hardware is believed to be located.

\paragraph{Assumed constants.}
\begin{itemize}
  \item NVIDIA H100 SXM5 TDP: 700\,W.
  \item NVIDIA A100 80GB TDP: 400\,W.
  \item NVIDIA L4 TDP: 72\,W.
  \item NVIDIA T4 TDP: 70\,W.
  \item Data-centre PUE: 1.1 (conservative; typical hyperscale PUE is
        1.10--1.15 \cite{google2023pue, aws2023pue}).
  \item US average grid intensity (EPA eGRID 2023): 0.367\,kg\,CO$_2$-eq / kWh
        \cite{epa2023egrid}.
  \item Indian average grid intensity (CEA 2023): 0.716\,kg\,CO$_2$-eq / kWh
        \cite{cea2023co2}.
  \item We use the US average for Modal H100 (US-East region as configured)
        and Tinker (location undisclosed; assumed US), and the Indian
        average for Colab Pro T4 used by authors based in Bengaluru
        (Google Cloud asia-south1).
\end{itemize}

\paragraph{GPU-hour estimates.}
Tinker GPU-hour counts are not disclosed by the platform. We estimate them
from wall-clock run duration and assume a single H100-class accelerator per
job, which is an unverified hardware-allocation assumption. The 18 Modal GPU-hours
in this scenario conflict with the 250 GPU-hours in Appendix~\ref{app:compute};
neither is selected as correct without a cohort-specific billing ledger.

\begin{table}[ht]
\centering
\small
\caption{Historical conditional carbon scenarios, not measured emissions. Tinker hardware
is undisclosed; 1$\times$ H100 per run is an assumption. Cohort boundaries and
failed-run coverage are unresolved; the scenario total is not a project estimate.}
\label{tab:carbon2}
\begin{tabular}{lrrrrr}
\toprule
\textbf{Platform} & \textbf{GPU-h} & \textbf{TDP\,(W)} & \textbf{PUE} &
\textbf{Energy\,(kWh)} & \textbf{CO$_2$\,(kg)} \\
\midrule
Tinker (assumed H100)   & \textasciitilde 950 & 700 & 1.1 & \textasciitilde 732 & \textasciitilde 269 \\
Modal H100 (US-East)    & \textasciitilde 18  & 700 & 1.1 & \textasciitilde 14  & \textasciitilde 5.1 \\
PES A100 (in-kind)      & \textasciitilde 60  & 400 & 1.1 & \textasciitilde 26  & \textasciitilde 19   \\
Colab Pro T4 (asia-south1) & \textasciitilde 40 & 70 & 1.2 & \textasciitilde 3.4 & \textasciitilde 2.4  \\
NVIDIA L4 (TRL baseline)   & \textasciitilde 10 & 72 & 1.1 & \textasciitilde 0.8 & \textasciitilde 0.3  \\
\midrule
\textbf{Scenario sum (not project total)} & & & & \textasciitilde 776 & \textasciitilde 296 \\
\bottomrule
\end{tabular}
\end{table}

\paragraph{Interpretation and uncertainty.}
The tabulated sum is conditional on unreconciled hours, assumed accelerator
counts, TDP, PUE, and locations. It is not a verified footprint or an upper bound:
unknown allocation can raise or lower actual consumption. Omitting failed jobs
in a future replication does not establish a bound on this project. A release
needs a single ledger with dates, cohort IDs, GPU counts, measured versus assumed
hours, credits versus cash, and failure status before reporting totals.

\paragraph{Offset and mitigation.}
We did not purchase carbon offsets. Instead, we have (a) made available selected
checkpoint artifacts so that downstream users do not
need to re-run our experiments; (b) retained local step-level CSV logs so that
learning curves can be inspected without re-training; and (c) documented
in Section~\ref{sec:compute_accounting} which runs were wasteful, so that
replicators can skip them. We regard reproducibility-by-artefact as a
more durable mitigation than offsets.

% ---------------------------------------------------------------------------
\subsection{Data Provenance}
\label{sec:provenance}
% ---------------------------------------------------------------------------

\ours{} uses only publicly released research datasets. No private data,
personally identifiable information, or licensed proprietary content is
used. No human annotators were employed during this work. We document
each dataset below.

\paragraph{GSM8K \cite{cobbe2021training}.}
Grade-school math word problems (7{,}473 train / 1{,}319 test)
authored by human writers contracted by OpenAI. Released under the
\textbf{MIT License} via \url{https://github.com/openai/grade-school-math}.
We use the standard train/test splits without modification. The canonical
citation is \citet{cobbe2021training}. Known limitations of GSM8K include
gender-neutral but culturally US-centric word problems (names, currencies,
sports) and a moderate rate of human-labelling errors estimated at
\textasciitilde 2\% by \cite{zhang2024gsm1k}. Our held-out evaluation
respects the test/train split.

\paragraph{Salesforce xLAM-Function-Calling-60k \cite{liu2024xlam}.}
60{,}000 function-calling examples generated by Salesforce Research as
training data for the xLAM agentic model family. Released under the
\textbf{CC BY 4.0} license via \url{https://huggingface.co/datasets/Salesforce/xlam-function-calling-60k}.
We use the public release as-is; specifically, we use the first 35 prompts
$\times$ 10 rollouts in the 10x Structural Ceiling tool-use track and the
full 60k split for the xlam-60k real-data run in Section~\ref{sec:results}.
Known limitations: xLAM schemas are synthetic and skew toward cleanly-typed
arguments; the dataset under-represents ambiguous, error-handling, and
multi-turn tool-calling. Attribution to Salesforce is required by the
license and is provided in Section~\ref{sec:acknowledgments} and in the
xLAM-derived model cards.

\paragraph{HumanEval \cite{chen2021codex}.}
164 Python programming problems with unit tests, released by OpenAI under
the \textbf{MIT License}. The local restricted sandbox does not implement an unmodified standard
pass@$k$ evaluation; code results are harness-bound reward probes. Known
limitations: small scale, English-language docstrings, single-language
coverage; documented contamination concerns against frontier pretraining
corpora \cite{riddell2024contamination}.

\paragraph{NuminaMath \cite{li2024numinamath}.}
Used by collaborator Mohammad Rafi for the multi-stage
GSM8K+NuminaMath pipeline. Released under \textbf{Apache 2.0}. We use a
downstream checkpoint reported by Rafi; our repository contains only his
scripts and the Apache-licensed derivative weights.

\paragraph{Open-Platypus \cite{lee2023platypus}.}
3{,}000 SFT examples used by collaborator Madhu for Qwen3-8B code
generation warm-up. Open-Platypus is a curated subset released under
\textbf{CC BY-NC 4.0}; the non-commercial restriction is respected in our
release, which is academic-use only.

\paragraph{Synthetic tool-use corpus (authored).}
We additionally generated a small five-tool synthetic corpus (calculator,
web-search stub, calendar, file-read, email-send) used for the reviewed descriptive
all-wrong observations. The corpus is authored by the paper's authors from
publicly documented APIs, contains no third-party content, and is
released under the \textbf{MIT License} in our repository under
\texttt{data/synthetic\_tools/}. Generation prompts are included for
auditability. No user data, scraped web content, or proprietary API
traces are present.

\paragraph{No web scraping, no human subjects.}
We did not scrape any website. We did not run any study with human
participants. No IRB review was therefore required. No data that could
reasonably be considered to contain PII, copyrighted content, or sensitive
personal attributes was used at any stage of training, evaluation, or
reward modelling.

% ---------------------------------------------------------------------------
\subsection{Closed-Source Tinker Acknowledgment}
\label{sec:tinker_closed}
% ---------------------------------------------------------------------------

A central tension in \ours{} is that a substantial fraction of our
headline numbers come from a closed-source commercial platform while our
paper simultaneously advocates for reproducibility and platform
independence. We do not resolve this tension by hiding it; we describe
it precisely.

\paragraph{What Tinker is and is not.}
Tinker \cite{thinkingmachines2024tinker} is a managed LLM fine-tuning and
inference service provided by Thinking Machines, Inc. It exposes a Python
SDK (\texttt{tinker==0.16.1}) that accepts custom loss functions
(\texttt{forward\_backward\_custom}), a limited set of optimisers, and
standard LoRA hyperparameters. Local scripts specify reward normalization, baseline subtraction, custom loss,
and requested LoRA/optimizer configuration. Undisclosed server-side details
include hardware placement, scheduling, execution numerics, and some batching
and sampler behavior. These are separate from the inspectable local objective.

\paragraph{What this means for our claims.}
We have not ruled out seed variance, checkpoint/model-size confounding, evaluator
differences, or backend effects. The arithmetic/GSM8K cross-library identity is
unresolved. Tinker results do not establish an upper bound or algorithmic
superiority. verl/OpenRLHF Task~4 dry runs are engineering artifacts, not
measured alternatives. The contribution remains the centered-reward identity,
descriptive audited cells, and reporting contract.

\paragraph{Reproducibility commitments we can make.}
Section~\ref{sec:reproducibility} distinguishes local source, observed
configuration, undisclosed internals, and evidence access. A dagger denotes only
interrupted execution; backend provenance is stated explicitly. Exact checkpoint
revisions and public access must be verified per artifact; no blanket release
or future rerun commitment is made here.

\paragraph{Key rotation and credential hygiene.}
Our Tinker API key (prefix \texttt{tml-...}) is held in a repository
\texttt{.env.example} placeholder only. The real key is stored in the
authors' password manager and rotated whenever a failure mode suggests
possible exfiltration. No Tinker key is committed to git history in
either the primary repository (\texttt{pes-llm-research/tinker-rl-lab}) or
the mirror (\texttt{arvindcr4/tinker-rl-lab}).

% ---------------------------------------------------------------------------
\subsection{Known Methodological Limits}
\label{sec:methodlimits}
% ---------------------------------------------------------------------------

In addition to the infrastructure failures and platform-specific caveats
enumerated in Section~\ref{sec:limitations}, we flag the following
methodological limits.

\paragraph{Short training horizons.}
Most reviewed Tinker runs have short, often 30-step horizons; this is not
a universal limit for all later scripts or campaigns. This is a
deliberate cost-control choice and means our Tinker numbers are
\emph{early-training snapshots}. We cannot rule out that longer training
would change the qualitative picture. The later 82\%$\to$83.3\% comparison
is one-sample fixed-baseline evidence, not paired-seed or reviewed-record
evidence; it is not a held-out gain whose size can be extrapolated.

\paragraph{Single-seed Tinker experiments.}
Each Tinker configuration was run once, not 3--5 times as best practice
prescribes \cite{henderson2018deep}. Tinker results therefore lack
variance estimates and do not support significance testing; we report
them descriptively. Seed-level inference requires the actual independent seed observations and
an identified estimand; reconstructed seed vectors are not eligible.

\paragraph{Train-set reward as primary Tinker metric.}
Tinker runs primarily report reward on training prompts. Only the GSM8K
track was followed up with a separate 200-example held-out evaluation
per seed. Other Tinker tracks (tool-use, xLAM) remain train-set-only and
we cannot distinguish memorisation from generalisation for those
settings.

\paragraph{LoRA only; no full fine-tuning.}
The reviewed LLM fine-tuning runs use LoRA \cite{hu2022lora}, with settings
varying by cohort; the Wave~6 rank sweep also includes rank 4. We have
\emph{not} tested whether full fine-tuning would re-order the
library/algorithm comparisons. The LoRA constraint is practical (cost) but
it means our claims about PPO vs.\ GRPO and TRL vs.\ Tinker are
LoRA-conditional.

\paragraph{Benchmark coverage is narrow.}
We evaluate four domains: GSM8K, MATH-500 (exploratory), HumanEval
(subset), and synthetic/xLAM tool-use. We do not evaluate on MT-Bench,
ArenaHard, safety benchmarks (HarmBench, ToxicChat), or truthfulness
benchmarks (TruthfulQA). Any claim about ``reasoning'' in the paper is
strictly a claim about the covered benchmarks.

\paragraph{No human preference data.}
We use verifiable rewards only (exact-match for math, unit-test for code,
schema-validity for tool-use). We do not study RLHF with human preference
data; findings should not be extrapolated to reward-model--based alignment
without further work.

\paragraph{Closed-source implementation opacity (reiterated).}
Comparisons between Tinker GRPO and TRL GRPO are confounded by the
closed-source nature of Tinker's implementation. See
Section~\ref{sec:tinker_closed} for detail.

\paragraph{Cross-platform hardware confounding.}
Tinker, Modal, Colab, and the PES A100 node use different accelerators,
different memory hierarchies, and different inference/training stacks
(Tinker proprietary, Modal vLLM, Colab PyTorch-native, TRL+Accelerate).
Observed differences in reward dynamics are not purely algorithmic.

\paragraph{Carbon estimates are approximate.}
As discussed in Section~\ref{sec:carbon}, Tinker GPU-hour and TDP are
inferred, not measured. Grid intensity is region-average, not marginal.
Numbers in Table~\ref{tab:carbon2} are conditional scenarios with
unreconciled cohorts, not measured emissions or upper bounds.

\paragraph{Exploratory, not confirmatory.}
This is an exploratory case study, not a pre-registered confirmatory
experiment. We did not pre-register primary hypotheses. The conflicting historical correction families and step-based significance
decisions are withdrawn (Section~\ref{sec:statistical-policy}); exploratory
coefficients do not establish calibrated inference.

\paragraph{Geographic and demographic limits.}
All authors are based at a single institution in Bengaluru, India. Our
choice of benchmarks, prompt phrasings, and evaluation design reflects
that context. Independent replication by groups in other regions, on
non-English tasks, and with non-Qwen / non-Llama families is an open
question.

% =============================================================================
% End of ethics_statement.tex
